\section{September 9, 2026}

\subsection{Matrix operations and linear maps}

Let $V$ and $W$ be vector spaces over $\bb{F}$ with ordered bases
\[
  \mathcal B=(v_1,\ldots,v_n)
  \qquad\text{and}\qquad
  \mathcal C=(w_1,\ldots,w_m).
\]
Suppose that $T,S\in\mathcal L(V,W)$ are represented by the matrices
$A=[A_{ij}]$ and $B=[B_{ij}]$, respectively. Thus
\[
  T(v_j)=\sum_{i=1}^m A_{ij}w_i
  \qquad\text{and}\qquad
  S(v_j)=\sum_{i=1}^m B_{ij}w_i.
\]
For each basis vector $v_j$,
\begin{align*}
  (T+S)(v_j)
  &=T(v_j)+S(v_j) \\
  &=\sum_{i=1}^m A_{ij}w_i+\sum_{i=1}^m B_{ij}w_i \\
  &=\sum_{i=1}^m(A_{ij}+B_{ij})w_i.
\end{align*}
Consequently, the matrix representing $T+S$ is obtained by adding the
matrices representing $T$ and $S$ entrywise.

\begin{definition}[Matrix addition]
  If $A,B\in M_{m\times n}(\bb{F})$, their sum is the matrix
  $A+B\in M_{m\times n}(\bb{F})$ defined by
  \[
    (A+B)_{ij}=A_{ij}+B_{ij}.
  \]
\end{definition}

Similarly, if $\alpha\in\bb{F}$, then
\[
  (\alpha T)(v_j)
  =\alpha T(v_j)
  =\sum_{i=1}^m \alpha A_{ij}w_i.
\]
Thus scalar multiplication of linear maps corresponds to entrywise
scalar multiplication of matrices.

\begin{definition}[Scalar multiplication of matrices]
  If $A\in M_{m\times n}(\bb{F})$ and $\alpha\in\bb{F}$, then
  $\alpha A\in M_{m\times n}(\bb{F})$ is defined by
  \[
    (\alpha A)_{ij}=\alpha A_{ij}.
  \]
\end{definition}

\subsection{Composition and matrix multiplication}

Let $Z$ be a third vector space with ordered basis
\[
  \mathcal D=(z_1,\ldots,z_p),
\]
and consider linear maps
\[
  T_2\colon V\longrightarrow W
  \qquad\text{and}\qquad
  T_1\colon W\longrightarrow Z.
\]
Suppose $T_2$ is represented by $B=[B_{kj}]\in
M_{m\times n}(\bb{F})$ and $T_1$ is represented by
$A=[A_{ik}]\in M_{p\times m}(\bb{F})$. For $j=1,\ldots,n$,
\begin{align*}
  (T_1\circ T_2)(v_j)
  &=T_1\left(\sum_{k=1}^m B_{kj}w_k\right) \\
  &=\sum_{k=1}^m B_{kj}T_1(w_k) \\
  &=\sum_{k=1}^m B_{kj}\sum_{i=1}^p A_{ik}z_i \\
  &=\sum_{i=1}^p\left(\sum_{k=1}^m A_{ik}B_{kj}\right)z_i.
\end{align*}
The coefficient of $z_i$ in this expression is the $(i,j)$ entry of
the matrix representing $T_1\circ T_2$.

\begin{definition}[Matrix multiplication]
  If $A\in M_{p\times m}(\bb{F})$ and
  $B\in M_{m\times n}(\bb{F})$, their product is the matrix
  $AB\in M_{p\times n}(\bb{F})$ whose entries are
  \[
    (AB)_{ij}=\sum_{k=1}^m A_{ik}B_{kj}.
  \]
\end{definition}

Thus composition of linear maps corresponds to multiplication of their
representing matrices, in the same order:
\[
  [T_1\circ T_2]=[T_1][T_2].
\]

Whenever all the products and sums involved are defined, matrix
multiplication satisfies
\begin{align*}
  (AB)C&=A(BC), \\
  A(B+C)&=AB+AC, \\
  (A+B)C&=AC+BC, \\
  A(\alpha B)&=\alpha(AB)=(\alpha A)B.
\end{align*}
In general, however, $AB\neq BA$; the two products may even have
different sizes, or one may be undefined. In particular,
$M_{n\times n}(\bb{F})$, with matrix addition, scalar multiplication,
and matrix multiplication, is an associative algebra over $\bb{F}$.

\subsection{Vector subspaces}

\begin{definition}[Subspace]
  Let $V$ be a vector space over $\bb{F}$. A subset $U\subseteq V$ is a
  \emph{subspace} of $V$ if $U$, with the addition and scalar
  multiplication inherited from $V$, is itself a vector space over
  $\bb{F}$.
\end{definition}

Checking all the vector-space axioms is unnecessary. The following
criterion reduces the verification to three conditions.

\begin{proposition}[Subspace criterion]\label{prop:subspace-criterion}
  A subset $U\subseteq V$ is a subspace of $V$ if and only if
  \begin{enumerate}[label=(\roman*)]
    \item $0\in U$;
    \item $u+v\in U$ whenever $u,v\in U$;
    \item $\alpha u\in U$ whenever $u\in U$ and $\alpha\in\bb{F}$.
  \end{enumerate}
  Equivalently, if $U$ is known to be nonempty, it is enough to check
  the two closure conditions (ii) and (iii).
\end{proposition}

\begin{example}
  The affine line
  \[
    H=\left\{\begin{bmatrix}x\\y\end{bmatrix}\in\bb{R}^2:x=2\right\}
  \]
  is not a subspace of $\bb{R}^2$, because it does not contain the zero
  vector.
\end{example}

The subspaces of $\bb{R}^2$ are precisely the zero subspace
$\{0\}$, the lines through the origin, and $\bb{R}^2$ itself.
